\cvsection{Expériences professionnelles}

\begin{cventries}
  \cventry
    {Stagiaire ingénieur de recherche en IA}
    {INRIA}
    {Bordeaux, France}
    {Mars 2026 -- Septembre 2026}
    {
      \begin{cvitems}
        \item {Conception et développement en autonomie, avec le suivi de mes encadrants, d’une application aidant les chercheurs à trouver des articles scientifiques à partir du besoin défini avec mon tuteur.}
        \item {Construction d’un agent multi-étapes sans framework d’orchestration avec intégration de l’API OpenAI, conception d’outils et gestion de l’état et de la mémoire.}
        \item {Conception de prompts adaptés au cas d’usage et sélection des messages par l’utilisateur pour contrôler le contexte transmis au LLM.}
        \item {Mise en œuvre d’un pipeline RAG avec préparation et découpage des documents, embeddings, indexation dans ChromaDB et recherche sémantique.}
        \item {Développement du backend Python/FastAPI avec API REST, OpenAPI et MongoDB, puis d’une interface React/Next.js destinée aux utilisateurs.}
        \item {Automatisation de la collecte de métadonnées scientifiques en ligne et transmission des résultats de tâches cron par webhooks.}
        \item {Définition d’un protocole de comparaison de LLM et de quantifications, analyse de la qualité et des performances, puis contribution à un article scientifique présentant les résultats et les limites.}
        \item {Conteneurisation de composants avec Docker et mise en place d’un pipeline CI/CD sur GitHub.}
      \end{cvitems}
    }

  \cventry
    {Stagiaire de recherche en apprentissage par renforcement}
    {Université Karunya}
    {Coimbatore, Inde}
    {Juin 2025 -- Septembre 2025}
    {
      \begin{cvitems}
        \item {Développement d’un modèle d’apprentissage par renforcement et de son environnement d’entraînement pour la navigation autonome d’un robot hospitalier.}
        \item {Déploiement et optimisation du logiciel et du modèle sur un Raspberry Pi afin de le tester en conditions réelles.}
      \end{cvitems}
    }
\end{cventries}
